\label{sec:results}

This section will be completed after the formal experiments are executed. The
planned result presentation has four parts.

\subsection{Baseline Results}

This subsection will report the public benchmark and railway-domain scores of
the ten-model pool before adaptation. The analysis should identify the strongest
general model, the strongest railway-domain baseline, and the smallest model
that remains competitive under the single-GPU constraint.

\begin{table*}[t]
\centering
\scriptsize
\caption{Completed baseline benchmark results on a single RTX 3090 24GB GPU. The precision column reports the best available public-benchmark run for each model in the current CSV export; Domain QA is shown for pipeline completeness and currently contains one exact-match test sample.}
\label{tab:main-results}
\begin{tabular}{llccccc}
\toprule
Model & Precision & MMLU & C-Eval & GSM8K & HumanEval & Domain QA \\
\midrule
Qwen3-0.6B & int8 & 0.401 & 0.434 & 0.000 & 0.183 & 0.000 \\
Qwen3-1.7B & int8 & 0.549 & 0.572 & 0.000 & 0.396 & 0.000 \\
Qwen3-4B & int8 & 0.681 & 0.715 & 0.000 & 0.652 & 0.000 \\
Qwen3-8B & int4 & 0.718 & 0.781 & 0.000 & 0.659 & 0.000 \\
Qwen2.5-7B-Instruct & bf16 & 0.718 & 0.796 & 0.000 & 0.665 & 0.000 \\
DeepSeek-R1-Distill-Qwen-7B & int8 & 0.527 & 0.492 & 0.000 & 0.116 & 0.000 \\
\bottomrule
\end{tabular}
\end{table*}


\subsection{Efficiency Results}

This subsection will compare bf16, int8, and int4 inference in terms of peak
VRAM, latency, and token throughput. The key question is whether a lower-bit
configuration preserves enough task performance to justify its memory and speed
benefits.

\begin{table*}[t]
\centering
\small
\caption{Deployment efficiency on the fixed efficiency prompt set. Lower latency and memory are better; higher generated-token throughput is better.}
\label{tab:efficiency-results}
\begin{tabular}{llccc}
\toprule
Model & Precision & Peak VRAM (GB) & Mean latency (s) & Tokens/s \\
\midrule
Qwen3-0.6B & bf16 & 1.15 & 4.72 & 54.28 \\
Qwen3-0.6B & int4 & 0.54 & 8.75 & 29.25 \\
Qwen3-0.6B & int8 & 0.77 & 21.20 & 12.08 \\
Qwen3-1.7B & bf16 & 3.24 & 4.84 & 52.88 \\
Qwen3-1.7B & int4 & 1.30 & 9.07 & 28.24 \\
Qwen3-1.7B & int8 & 1.95 & 21.36 & 11.99 \\
Qwen3-4B & bf16 & 7.54 & 6.61 & 38.74 \\
Qwen3-4B & int4 & 2.55 & 11.44 & 22.38 \\
Qwen3-4B & int8 & 4.17 & 25.48 & 10.05 \\
Qwen3-8B & bf16 & 15.30 & 6.45 & 39.69 \\
Qwen3-8B & int4 & 5.78 & 11.40 & 22.46 \\
Qwen3-8B & int8 & 8.85 & 25.13 & 10.19 \\
Qwen2.5-7B-Instruct & bf16 & 14.21 & 5.49 & 46.60 \\
Qwen2.5-7B-Instruct & int4 & 13.64 & 8.15 & 31.41 \\
Qwen2.5-7B-Instruct & int8 & 8.16 & 20.15 & 12.71 \\
DeepSeek-R1-Distill-Qwen-7B & bf16 & 14.21 & 5.49 & 46.64 \\
DeepSeek-R1-Distill-Qwen-7B & int4 & 13.64 & 8.23 & 31.12 \\
DeepSeek-R1-Distill-Qwen-7B & int8 & 8.26 & 19.30 & 13.26 \\
\bottomrule
\end{tabular}
\end{table*}


\subsection{QLoRA Adaptation Results}

This subsection will compare railway-domain performance before and after QLoRA
for the selected adaptation candidates. It should also report training cost,
adapter size, and whether adaptation causes any degradation on public tasks.

\begin{table*}[t]
\centering
\scriptsize
\caption{Current QLoRA adapter evaluation export. The Domain QA exact-match score is not yet a mature domain benchmark result because the current adapter evaluation summary contains one test sample.}
\label{tab:qlora-results}
\begin{tabular}{lccccc}
\toprule
Model & Baseline Domain & Adapter Domain & Adapter MMLU & Adapter C-Eval & Tokens/s \\
\midrule
Qwen3-4B & 0.000 & 0.000 & -- & -- & 13.149 \\
Qwen2.5-7B-Instruct & 0.000 & 0.000 & 0.710 & 0.789 & 18.032 \\
\bottomrule
\end{tabular}
\end{table*}


\subsection{Model Selection Analysis}

This subsection will synthesize the results into deployment guidance. The final
recommendation should distinguish between the best accuracy-oriented model, the
best efficiency-oriented model, and the best model after railway-domain
adaptation.
